\documentclass[tikz,border=8pt]{standalone}
\usepackage{tikz}
\usepackage{amsmath}
\usepackage{amssymb}
\usepackage{pgfplots}
\pgfplotsset{compat=1.18}
\usepackage{ctex}
\usetikzlibrary{calc,positioning,arrows.meta,matrix,trees}

% Union-Find 路径压缩演示
% 初始：5节点链状树 0-1-2-3-4（parent[i]=i-1, root=0）
% 展示 find(4) 路径压缩过程 + parent/rank 数组演化

\begin{document}
\begin{tikzpicture}[
  every node/.style={font=\small},
  treenode/.style={draw=blue!60, circle, minimum size=0.75cm, thick, fill=blue!10},
  rootnode/.style={draw=orange!70!black, circle, minimum size=0.75cm, very thick, fill=orange!15},
  compnode/.style={draw=green!60!black, circle, minimum size=0.75cm, thick, fill=green!15},
  cell/.style={draw=gray!60, rectangle, minimum width=0.75cm, minimum height=0.6cm, thick, fill=white},
  hcell/.style={draw=orange!70, rectangle, minimum width=0.75cm, minimum height=0.6cm, thick, fill=orange!10},
  gcell/.style={draw=green!60!black, rectangle, minimum width=0.75cm, minimum height=0.6cm, thick, fill=green!10},
  ptr/.style={->,>=Stealth,thick},
]

%% 标题
\node[font=\large\bfseries] at (6.0, 1.5)
  {Union-Find 路径压缩（Path Compression）};
\node[font=\small,gray] at (6.0, 0.9)
  {5 节点链状树，find(4) 后所有节点直接指向根节点 0};

%% ─── 阶段 A：初始状态 ─────────────────────────────────────────
\node[font=\small\bfseries,blue!70] at (2.2, 0.3) {阶段 A：初始链状树};

% 树节点（竖链）: 0 是根，1→0，2→1，3→2，4→3
\node[rootnode] (N0) at (2.2, -0.5) {0};
\node[treenode] (N1) at (2.2, -1.5) {1};
\node[treenode] (N2) at (2.2, -2.5) {2};
\node[treenode] (N3) at (2.2, -3.5) {3};
\node[treenode] (N4) at (2.2, -4.5) {4};

\draw[ptr] (N1) -- (N0);
\draw[ptr] (N2) -- (N1);
\draw[ptr] (N3) -- (N2);
\draw[ptr] (N4) -- (N3);

% parent 数组（A 阶段）
\node[font=\scriptsize,gray] at (0.5, -0.5) {parent:};
\node[font=\tiny,gray] at (0.5, -5.1) {（初始）};

\foreach \i/\p in {0/0,1/0,2/1,3/2,4/3} {
  \node[cell,font=\scriptsize] at (0.0+\i*0.77, -5.5) {\p};
  \node[font=\tiny,gray] at (0.0+\i*0.77, -5.9) {\i};
}
\node[font=\tiny,gray,anchor=west] at (3.9, -5.5) {← parent[]};

%% ─── 阶段 B：find(4) 递归追溯路径 ──────────────────────────────
\node[font=\small\bfseries,blue!70] at (7.0, 0.3) {阶段 B：find(4) 递归查根};

\node[rootnode] (M0) at (7.0, -0.5) {0};
\node[treenode] (M1) at (7.0, -1.5) {1};
\node[treenode] (M2) at (7.0, -2.5) {2};
\node[treenode] (M3) at (7.0, -3.5) {3};
\node[compnode] (M4) at (7.0, -4.5) {4};

\draw[ptr] (M1) -- (M0);
\draw[ptr] (M2) -- (M1);
\draw[ptr] (M3) -- (M2);
\draw[ptr,very thick,green!60!black] (M4) -- (M3);

% 路径追溯标注（右侧）
\node[font=\tiny,green!60!black,anchor=west] at (7.45, -4.5) {← 从此节点开始 find};
\node[font=\tiny,orange!70!black,anchor=west] at (7.45, -0.5) {← 根节点};

%% ─── 阶段 C：路径压缩后 ─────────────────────────────────────────
\node[font=\small\bfseries,blue!70] at (11.5, 0.3) {阶段 C：路径压缩后};

\node[rootnode] (P0) at (11.5, -0.5) {0};
\node[compnode] (P1) at (10.0, -1.8) {1};
\node[compnode] (P2) at (11.0, -1.8) {2};
\node[compnode] (P3) at (12.0, -1.8) {3};
\node[compnode] (P4) at (13.0, -1.8) {4};

\draw[ptr,green!60!black,thick] (P1) -- (P0);
\draw[ptr,green!60!black,thick] (P2) -- (P0);
\draw[ptr,green!60!black,thick] (P3) -- (P0);
\draw[ptr,green!60!black,thick] (P4) -- (P0);

\node[font=\tiny,green!60!black] at (11.5, -2.3)
  {所有节点直接指向根 0};

% parent 数组（C 阶段）
\node[font=\tiny,gray] at (9.5, -5.1) {（压缩后）};
\foreach \i/\p in {0/0,1/0,2/0,3/0,4/0} {
  \node[gcell,font=\scriptsize] at (9.0+\i*0.77, -5.5) {\p};
  \node[font=\tiny,gray] at (9.0+\i*0.77, -5.9) {\i};
}
\node[font=\tiny,gray,anchor=west] at (12.9, -5.5) {← parent[]};

%% rank 数组说明框
\node[draw=gray!40,fill=white,thin,rounded corners=2pt,
      font=\scriptsize,inner sep=6pt,align=left]
  at (6.0, -7.0)
  {\textbf{rank 数组}（按秩合并辅助）：\\[3pt]
   rank[] = [0, 0, 0, 0, 0]（初始所有节点秩为 0）\\[2pt]
   union(a,b)：rank[a]<rank[b] $\Rightarrow$ parent[a]=b\\[2pt]
   rank[a]==rank[b] $\Rightarrow$ parent[b]=a, rank[a]++};

%% 算法说明
\node[draw=gray!50,fill=white,rounded corners=3pt,font=\small,
      inner sep=8pt,align=left]
  at (6.0, -9.0)
  {\textbf{Union-Find 路径压缩 + 按秩合并（近乎 O(1) 均摊）：}\\[4pt]
   \texttt{find(x)}: 若 parent[x]!=x，递归找根并令 parent[x]=根（路径压缩）\\[2pt]
   \texttt{union(a,b)}: 找到各自根，按 rank 合并（低秩接到高秩下）\\[2pt]
   两者配合，每次操作均摊近乎 O($\alpha$(n))（逆 Ackermann 函数）};

\end{tikzpicture}
\end{document}
